\section{Proofs}\label{sectionproofs}


\subsection[Tensor $\J$]{Tensor $\boldsymbol{\J}$}\label{sectionJ}

The multiplication map $\mult$ in $\Z_n$ (we return to the original notation) is given by the
prescription~\rf{not5a}, as in any reduction algebra. It can be formally expanded into
a series over the root lattice of certain bilinear maps as follows. Set
\begin{gather*}
\UU(\bb_\pm):=\Uh\otimes_{\U(\h)}\U(\bb_\pm)  , \qquad \UUb:=\UU(\bb_-)\otimes_{\Uh}  \UU(\bb_+)  .
\end{gather*}
All these are associative algebras. Besides, both algebras~$\UU(\bb_\pm)$ are $\Uh$-bimodules. The
algebra~$\UUb$ admits three commuting actions of~$\Uh$. Two of them are given by the assignments
\begin{gather*}
X(Y\otimes Z):=XY\otimes Z  ,\qquad (Y\otimes Z)X:=Y\otimes ZX  ,
\end{gather*}
for any $X\in\Uh$, $Y\in\UU(\bb_-)$ and $Z\in\UU(\bb_+)$.
The third action associates to any $X\in\Uh$, $Y\otimes Z\in\UUb$ the element $YX\otimes Z=Y\otimes XZ\in\UUb$.


Present the projector $P$ in an ordered form:
\begin{gather} \label{proof1}
P=
\sum_{\gamma,i}\grave{F}_{\gamma,i}\grave{E}_{\gamma,i}\grave{H}_{\gamma,i}=
\sum_{\gamma,i}\grave{H}_{\gamma,i}\grave{F}_{\gamma,i}\grave{E}_{\gamma,i}  ,
\end{gather}
the summation is over $\gamma\in\Q_+$ and $i\in {\mathbb{Z}}_{\geq 0}$;
every $\grave{F}_{\gamma,i}$ is an element of $\U(\n_-)$ of the weight~$-\gamma$, every
$\grave{E}_{\gamma,i}$ is an element of $\U(\n_+)$ of the weight $\gamma$ and
$\grave{H}_{\gamma,i}\in\Uh$. Let $\J$ be the following element of $\UUb$:
\begin{gather*}
\J :=\sum_{\gamma,i}\grave{F}_{\gamma,i}\otimes\grave{E}_{\gamma,i}\grave{H}_{\gamma,i}=
\sum_{\gamma,i}\grave{H}_{\gamma,i}\grave{F}_{\gamma,i}\otimes\grave{E}_{\gamma,i}  ,\qquad \gamma\in\Q_+  ,\quad i\in {\mathbb{Z}}_{\geq 0}  .
\end{gather*}
Due to the PBW theorem in $\U(\gl_n)$ the tensor~$\J$ is uniquely def\/ined by the projector $P$; it is
of total weight zero: $h\J=\J h$ for any $h\in\h$. We have the weight decomposition of $\J$ with
respect to the adjoint action of $\h$ in the second tensor factor of~$\UUb$:
\begin{gather*}
\J= \bigoplus_{\lambda\in\Q_+}\J_\lambda  ,
\end{gather*}
where $\J_\lambda$ consists of all the terms, corresponding to $\grave{F}_{\lambda,i}\grave{E}_{\lambda,i}
\grave{H}_{\lambda,i}$ in~\rf{proof1}
(contributing to $\lambda\in\Q_+$ in the summation),
\begin{gather*}
\J_\lambda :=\sum_i\grave{F}_{\lambda,i}\otimes\grave{E}_{\lambda,i}\grave{H}_{\lambda,i}  .
\end{gather*}
By def\/inition of $\J$, the multiplication $\mult$ in the double coset space $\Z_n$ can be described by the relation
\begin{gather}\label{proof4}
a\mult b= m\left( (a\otimes 1)\J (1\otimes b )\right),
\end{gather}
where $m(\sum_ic_i\otimes d_i)$ is the image in $\Z_n$ of the element $\sum_ic_id_i$.
Moreover, in \rf{proof4} we can replace all products $\grave{E}_{\gamma,i}b$ in the second tensor
factor by the adjoint action of $\grave{E}_{\gamma,i}$ on $b$ (in fact, for $\grave{E}_{\gamma,i}=e_{\gamma_m}\cdots e_{\gamma_1}$,
we can replace $\grave{E}_{\gamma,i}b$ by $[\grave{E}_{\gamma,i},b]$ or by $\hat{e}_{\gamma_m}\cdots \hat{e}_{\gamma_1}(b)$,
see \rf{dehaco}) and likewise all products $a\grave{F}_{\gamma,i}$
in the f\/irst tensor factor by the opposite adjoint action of $\grave{F}_{\gamma,i}$ on $a$.
We have a~decomposition of the product $\mult$ into a sum over $\Q_+$:
\begin{gather}\label{proof5}
a\mult b=\sum_{\lambda\in\Q_+}a\mult_\lambda b  ,\qquad\text{where}\quad
a\mult_\lambda b:=m\left( (a\otimes 1)  \J_\lambda (1\otimes b )\right)  .
\end{gather}
If $a$ and $b$ are weight elements of $\Z_n$ of weights $\nu(a)$ and $\nu(b)$, then the product
$a\mult_\lambda b$ is the image in $\Z_n$ of the sum $\sum_i a_ib_i$, where the weight of each
$b_i$ is $\nu(b)+\lambda$, and the weight of each~$a_i$ is $\nu(a)-\lambda$.


The tensor $\J$ satisf\/ies the Arnaudon--Buf\/fenoir--Ragoucy--Roche (ABRR) dif\/ference equation~\cite{ABRR}, see also~\cite{K} for the translation of the results of~\cite{ABRR} to the language
of reduction algebras. To describe the equation, let  $\vartheta=\frac{1}{2}
\sum_{k=1}^n\th_k^2\in \U(\h)$;
for any positive root $\gamma\in\Delta_+$, denote by~$T_\gamma$ the following linear operator on the vector space~$\UUb$:
\begin{gather*}T_{\gamma}(X\otimes Y):=X e_{-\gamma}\otimes e_{\gamma}Y  .\end{gather*}
The ABRR equation means the relation \cite{ABRR,K}:
\begin{gather*}%\label{proof2}
[1\otimes \vartheta, \J]=-\sum_{\gamma\in\Delta_+}T_{\gamma}(\J)
 .
\end{gather*}
This relation is equivalent to the following system of recurrence relations for the weight components
$\J_\lambda$:
\begin{gather} \label{proof3}
\J_\lambda\cdot\left(\th_\lambda+\frac{(\lambda,\lambda)}{2}\right)=
-\sum_{\gamma\in\Delta_+}T_\gamma\left(\J_{\lambda-\gamma}\right)  ,
\end{gather}
where $\th_\lambda :=\sum_k\lambda_k\th_k$ for $\lambda=\sum_k\lambda_k\ve_k$. The recurrence
relations~(\ref{proof3})  together with the initial condition $\J_0=1\otimes 1$
uniquely determine all weight components~$\J_\lambda$.


It should be noted that the recurrence relations \rf{proof3} provides less information about the structure of the denominators (from~$\U (\h )$)
of the summands of the extremal projector $P$ than the information implied
 by the product formula (see~\cite{AST}) for the extremal projector.


Using \rf{proof3} we get in particular:
\begin{gather} \label{proof3a}
\J_{\a}  = -(\th_{\a}+1)^{-1}e_{-a}\otimes e_{\a} ,
\qquad \a=\ve_i-\ve_{i+1} ,
\\ \notag
 \J_{\a+\beta}  =  (\th_{\a+\b}+1)^{-1}\Bigl(-e_{-\a-\b}\otimes e_{\a+\b}+
(\th_{\a}+1)^{-1} e_{-\a}e_{-\b}\otimes e_{\b}e_{\a}\\
\phantom{\J_{\a+\beta}  =}{}  +(\th_{\b}+1)^{-1} e_{-\b}e_{-\a}\otimes e_{\a}e_{\b}\Bigr)
  ,\qquad \a=\ve_{i-1}-\ve_{i}  ,\quad \b=\ve_{i}-\ve_{i+1}  ,\label{proof3b}
\\
  \label{proof3c}\J_{\ve_i-\ve_{j}+\ve_k-\ve_{l}}= \J_{\ve_i-\ve_{j}}\cdot
\J_{\ve_k-\ve_{l}}  ,\qquad  i<j<k<l  .
\end{gather}

\subsection{Braid group action}

%{\bf 1.}
The proof of the relations \rf{acwot} and \rf{3.5} consists of the following arguments, valid for any
reduction algebra. Let $\a$ be any simple root of $\gl_n$, $\a=\ve_i-\ve_{i+1}$ and $\g_\a$ the
corresponding $\sl_2$ subalgebra of $\gl_n$. It is spanned by the elements $e_\a=e_{i,i+1}$,
$e_{-\a}= e_{i+1,i}$ and $h_\a= h_i-h_{i+1}$. Let $\sy_\a=\sy_i$ be the corresponding
automorphism of the algebra $\Ar$ and $\q_\a=\q_i$ the Zhelobenko automorphism of
$\Z_n$. Assume that $Y\in\Ar$ belongs, with respect to the adjoint action of $\g_\a$, to an irreducible
f\/inite-dimensional $\g_\a$-module of dimension $2j+1$, $j\in\{\ts0,1/2,1,\ldots\ts\}$. Assume further
that $Y$ is homogeneous, of weight $2\ts m$, $[h_\a,Y]=2m Y$. Identify $Y$ with its image in $\Z_n$. Then
$\q_\a(Y)$ coincides with the image in $\Z_n$ of the element
\begin{gather*}%\label{proof6}
\prod_{i=1}^{j+m}\,(h_\a+i+1)\ts\cdot\ts\sy_\alpha(Y)\cdot
{\ds\prod_{i=1}^{j+m} (h_\a-i+1)^{-1}}  .
\end{gather*}
This can be checked directly using \cite[Proposition 6.5]{KO}.


In the realization of irreducible $\sl_2$-modules as the spaces of homogeneous polynomials in two
variables $u$ and $v$,
\begin{gather*}e_\a\mapsto u\frac{\partial}{\partial v}  , \qquad h_\a\mapsto u\frac{\partial}{\partial u}-v\frac{\partial}{\partial v}
\qquad {\mathrm{and}}\qquad  e_{-\a}\mapsto v\frac{\partial}{\partial u}  ,
\end{gather*}
the operator $\sy_\alpha$ becomes $(\sy_\alpha f)(u,v)=f(-v,u)$, or, in the basis $|j,k\rangle :=x^{j+k}y^{j-k}$ ($j$ labels the
representation; $k=0,1,\dots ,2j$),
\begin{gather*}\sy_\alpha : \  |j,-j+k\rangle\mapsto (-1)^k |j,j-k\rangle  .\end{gather*}

\begin{proof}[Proof  of Lemma \ref{lqw02}, Subsection~\ref{brgac}.] %, Subsection \ref{brgac}.
To see this, write a reduced expression for $\q_{w_0}$, $\q_{w_0}=\q_{\alpha_{i_1}}\cdots \q_{\alpha_{i_M}}$
with $\alpha_{i_1},\dots ,\alpha_{i_M}$ simple roots. Then $\q_{w_0}=\q_{\alpha_{i_M}}\cdots \q_{\alpha_{i_1}}$ as well.
Writing, for $\q_{w_0}^2$, the second expression after the f\/irst one, we get squares of
$\q_{\alpha_{i_s}}$'s (which are conjugations by $\th_{\alpha_{i_s}}^{-1}$'s; they thus commute) one after another.
Moving these conjugations to the left through the remaining~$\q$'s, we produce, exactly like in the
construction of a system of all positive roots from a reduced expression for
the longest element of the Weyl group of a reductive Lie group, the conjugation by the product (\ref{conjq02})
over all positive roots.
\end{proof}

\begin{proof}[Proof  of Proposition~\ref{paqw0}, Subsection~\ref{brgac}.] %, subsection \ref{brgac}.
Only formula \rf{3.6} needs a proof (formula~\rf{3.6a}
is a particular case of~\rf{3.2}).


For a moment, denote the longest element
of the symmetric group $S_n$ by $\q_{w_0}^{(n)}$. Let $\psi_j:=\q_j\q_{j-1}\cdots\q_1$ (the product in the descending order).
We have $\q_{w_0}^{(n+1)}=\q_{w_0}^{(n)}\psi_n$ and  $\q_{w_0}^{(n+1)}=\psi_1\psi_2\cdots\psi_n$ (the product in the
ascending order).


For $j<n$ it follows from \rf{3.5} that $\psi_j (\s_{n+1,1})=(-1)^j\s_{n+1,j+1}$ (say, by induction on $j$). So,
\begin{gather*}\psi_n(\s_{n+1,1})=q_n\psi_{n-1} (\s_{n+1,j+1})=(-1)^{n-1}q_n(\s_{n+1,n})=(-1)^{n}\s_{n,n+1}  ,\end{gather*}
again by \rf{3.5}. Next, $\psi_k\psi_{k+1}\cdots\psi_{n-1}(\s_{n,n+1})=\s_{k,n+1}$ by induction on $n-k$ and
again \rf{3.5}. Thus,
\begin{gather}\label{binh}
\q_{w_0}(\s_{n+1,1})=(-1)^n \s_{1,n+1} ,
\end{gather}
establishing \rf{3.6} for $i=n+1$ and $j=1$. We now prove \rf{3.6} for $i>j$ (positions below the main diagonal) by induction backwards
on the height $i-j$ of a negative root;  the formula~\rf{binh} serves as the induction base. Assume that~\rf{3.6}
is verif\/ied for a given level $i-j$ and $i-j-1>0$ (so that the positions $(i,j+1)$ and $(i-1,j)$ are still under the main
diagonal). By \rf{3.5}, $\s_{i,j+1}=-\q_j(\s_{ij})$, therefore
\begin{gather*}\q_{w_0}(\s_{i,j+1}) = -\q_{w_0}(\q_j(\s_{ij}))=-\q_{j'-1}(\q_{w_0}(\s_{ij})) \\
\phantom{\q_{w_0}(\s_{i,j+1})}{}  = (-1)^{i+j+1} \q_{j'-1}\left(\s_{i'j'}\ds{\prod_{a:a<i'}}\tphi_{ai'}  \ds{\prod_{b:b>j'}}\tphi_{j'b}\right)\\
\phantom{\q_{w_0}(\s_{i,j+1})}{} = (-1)^{i+j+1}\s_{i',j'-1}\tphi_{j'-1,j'}\ds{\prod_{a:a<i'}}\tphi_{ai'}  \ds{\prod_{b:b>j'}}\tphi_{j'-1,b}\\
\phantom{\q_{w_0}(\s_{i,j+1})}{} = (-1)^{i+j+1}\s_{i',(j+1)'}\ds{\prod_{a:a<i'}}\tphi_{ai'}\ds{\prod_{b:b>(j+1)'}}\tphi_{(j+1)',b}  .
\end{gather*}
In the second equality we used the identity $\q_{w_0}\q_j=\q_{j'-1}\q_{w_0}$ in the braid group; the third equality is the
induction assumption; in the fourth equality we used that $i'\neq j'-1$ (since $i-j-1>0$) and then~\rf{3.5}; in the f\/ifth equality we
replaced $j'-1$ by $(j+1)'$. The calculation for $\q_{w_0}(\s_{i-1,j})$ is similar; it uses $\s_{i-1,j}=\q_{i-1}(\s_{ij})$.
The proof of the formula~\rf{3.6} for positions below the main diagonal is f\/inished.


The proof of  \rf{3.6} for $i<j$ (positions above the main diagonal) follows now from Lem\-ma~\ref{lqw02}. %, Subsection \ref{brgac}.
\end{proof}

\subsection{Derivation of relations}\label{sectionDerRel}

The set  of def\/ining relations in $\Z_n$ divides into several dif\/ferent types,
see Section~\ref{section3.3}. We prove the necessary amount of relations
of each type and get the rest by applying the transformations from the braid group as well as the anti-involution~$\epsilon$, see~(\ref{anep}).


We never use the automorphism $\omega$, def\/ined in \rf{not2a}, in the derivation of relations. However, the involution $\omega$ is compatible with our set
of relations in the sense explained in Section \ref{proofmain}.


In the following we denote by the symbol $ \equiv $ the equalities of elements from $\Ab$ modulo the sum $(\Jb+\Ib )$ of two ideals $\Jb$ and $\Ib$
def\/ined in the beginning of Section \ref{section2}. Moreover,
for any two elements $X$ and $Y$ of the algebra $\Ab$ we may regard the expressions $X\mult Y$ and $X\mult_\lambda Y$ as the sums
of elements from $\Ab$ def\/ined in \rf{proof4} and \rf{proof5}. The sum $X\mult_\lambda Y$  is f\/inite. By the construction,
all but a f\/inite number of terms in the product $X\mult Y$ belong to $(\Jb+\Ib )$. Unlike to the system of notation
adopted in Section~\ref{section2}, our proof of each relation in $\Z_n$ will use equalities in $\Ab$ taken modulo $(\Jb+\Ib )$.


We also use the notation $H_i$ for the element $E_{ii}\in\Ar$ and $H_{ij}=H_i-H_j=E_{ii}-E_{jj}$.

{\bf 1.} We f\/irst prove in $\Z_n$ the relation
\begin{gather}\label{proof7}
\s_{12}\mult\s_{13}=\s_{13}\mult\s_{12}  \frac{\th_{23}}{\th_{23}+1}  .
\end{gather}
Elements $\s_{12}$ and $\s_{13}$ are images in $\Z_n$ of $E_{12}$ and $E_{13}$. Consider the
product $E_{12}\mult_\lambda E_{13}$. Since the adjoint action of $\gl_n$ preserves the space
$\p$, see Section~\ref{section-notation}, this product is the sum of such monomials $E_{ij}E_{kl}$,
with coef\/f\/icients in $\Uh$, that (i): the weight $\ve_k-\ve_l$
of $E_{kl}$ is equal to the weight $\ve_{1}-\ve_3$ of $E_{13}$ plus $\lambda\in\Q_+$, while (ii): the
weight $\ve_i-\ve_j$ of $E_{ij}$ is equal to the weight $\ve_{1}-\ve_2$ of $E_{12}$ minus
$\lambda$. Assume that $E_{12}\mult_\lambda E_{13}\neq 0$. By (i), $\lambda =-\ve_1+\ve_3+\ve_k-\ve_l$
and it can be positive only if $k=1$ and $l\geq 3$. So, the condition (i) implies that either
$\lambda=0$ or $\lambda=\ve_3-\ve_l$ with $l>3$. The possibility $\lambda=\ve_3-\ve_l$, $l>3$, is excluded
by the condition (ii). Therefore, $\lambda=0$ and
\begin{gather}
\label{proof8}
E_{12}\mult E_{13}\equiv E_{12}E_{13} .
\end{gather}
Similarly, for $\lambda\in\Q_+$, which can non-trivially contribute to the product $E_{13}\mult E_{12}$,
the analogue of the condition (i) on the weight $\lambda$ gives the
restriction $\lambda =0$ or $\lambda=\ve_2-\ve_k$, $k>2$; the analogue of the condition (ii)
further restricts $\lambda$:  $\lambda=0$ or $\lambda=\ve_2-\ve_3$, so we have
\begin{gather*}
E_{13}\mult_{\ve_2-\ve_3} E_{12}\equiv -E_{13}e_{32}e_{23}  \frac{1}{\th_{23}+1}E_{12}\equiv
-E_{13}e_{32}e_{23}E_{12}  \frac{1}{\th_{23}}\equiv E_{12}E_{13}  \frac{1}{\th_{23}}  ,
\end{gather*}
since $\J_{\ve_2-\ve_3}=-e_{32}\otimes e_{23}(\th_{23}+1)^{-1}$ as it follows from the ABRR
equation, see~(\ref{proof3a}), or from the precise explicit expression for the projector~$P$,
see~\cite{AST}. Thus, since $E_{12}$ and $E_{13}$ commute in the universal enveloping algebra
\begin{gather}\label{proof9}E_{13}\mult E_{12}\equiv E_{13}E_{12}+E_{13}\mult_{\ve_2-\ve_3} E_{12}=
E_{12}E_{13}\left(1+\frac{1}{\th_{23}}\right),
\end{gather}
Comparing \rf{proof8} and \rf{proof9} we f\/ind \rf{proof7}.


Applying to \rf{proof7} the anti-involution $\epsilon$, see (\ref{anep}), we get the relation
\begin{gather}\label{proof10}
\s_{21}\mult\s_{31}=\s_{31}\mult\s_{21}  \frac{\th_{23}+1}{\th_{23}}  .
\end{gather}

The rest of the relations \rf{relation1} are obtained from~\rf{proof7} and~\rf{proof10}
by applying dif\/ferent transformations $\q_w$ from the Weyl group.

{\bf 2.} Now we prove in $\Z_n$ the relation
\begin{gather}\label{proof11}
\s_{13}\mult\s_{24}-\s_{24}\mult\s_{13}=
\left(\frac{1}{\th_{12}}-\frac{1}{\th_{34}}\right)\s_{23}\mult\s_{14}  .
\end{gather}
We begin by the proof of this relation in $\Z_4$. We proceed in the same manner as for the derivation of the
relation (\ref{proof7}),
\begin{gather*}E_{13}\mult E_{24}\equiv E_{13}E_{24}+E_{13}\mult_{\ve_1-\ve_2} E_{24} \\
\phantom{E_{13}\mult E_{24}}{} \equiv
E_{13}E_{24}-E_{13}e_{21}e_{12}  \frac{1}{\th_{12}+1} E_{24}
 \equiv E_{13}E_{24}+E_{23}E_{14} \frac{1}{\th_{12}}  ,\\
 E_{24}\mult E_{13}\equiv E_{24}E_{13}+E_{24}\mult_{\ve_3-\ve_4} E_{13}\\
 \phantom{E_{24}\mult E_{13}}{}  \equiv
E_{24}E_{13}-E_{24}e_{43}e_{34}  \frac{1}{\th_{34}+1} E_{13} \equiv E_{13}E_{24}+E_{23}E_{14}  \frac{1}{\th_{34}}  ,\\
 E_{23}\mult E_{14} \equiv E_{23}E_{14}  .
 \end{gather*}
Combining the three latter equalities we obtain \rf{proof11} in $\Z_4$.


The dif\/ference of the left and right hand sides of \rf{proof11} in $\Z_n$ is a linear combination of monomials in $\s_{ij}$ of
the total weight $\ve_1+\ve_2-\ve_3-\ve_4$. The weight is non-trivial, so the monomials can be only quadratic.
Due to the stabilization phenomenon, each monomial should contain $\s_{ij}$ with $i>4$ or $j>4$, but, by the weight arguments,
there is no such non-zero possibility, which completes the proof of the relation \rf{proof11} in $\Z_n$.


The rest of relations \rf{relation2} is then obtained by applications of the transformations from the braid group.

{\bf 3a.} We continue and derive in $\Z_4$ the relation (we remind the notation $t_{ij}:=\s_{ii}-\s_{jj}$,
see Section~\ref{section2},
and $H_{ij}=E_{ii}-E_{jj}$):
\begin{gather} \label{proof12}
\s_{23}\mult\s_{12}-\s_{12}\mult\s_{23}=t_{12}\mult\s_{13}
 \frac{1}{\th_{12}}+t_{23}\mult\s_{13}  \frac{1}{\th_{23}}-\s_{43}\mult\s_{14}  \frac{\th_{34}+1}{\th_{34}\th_{24}}  .
\end{gather}
Using \rf{proof3a}--\rf{proof3c}, we calculate, to obtain the result for $\Z_4$:
\begin{gather}E_{12}\mult E_{23} \equiv E_{12}E_{23}+ E_{12}\mult_{\ve_1-\ve_2} E_{23}+
E_{12}\mult_{\ve_1-\ve_2+\ve_3-\ve_4} E_{23} \nonumber\\
\phantom{E_{12}\mult E_{23}}{}
 \equiv E_{12}E_{23}-H_{12}E_{13} \frac{1}{\th_{12}}  ,\label{1223r1}\\
  E_{23}\mult E_{12} \equiv E_{23}E_{12}+ E_{23}\mult_{\ve_2-\ve_3} E_{12}+
E_{23}\mult_{\ve_2-\ve_4} E_{12}\notag\\
\phantom{E_{23}\mult E_{12}}{} \equiv E_{23}E_{12}+H_{23}E_{13}  \frac{1}{\th_{23}}-E_{43}E_{14} \frac{(\th_{23}-1)}{\th_{23}
\th_{24}} ,\label{1223r2}\\
  H_{12}\mult E_{13} \equiv H_{12} E_{13}+H_{12}\mult_{\ve_3-\ve_4} E_{13}\equiv
H_{12} E_{13}  ,\label{1223r3}\\
  H_{23}\mult E_{13} \equiv H_{23} E_{13}+H_{23}\mult_{\ve_3-\ve_4} E_{13}\equiv
H_{23} E_{13}+E_{43}E_{14}  \frac{1}{\th_{34}}  ,\label{1223r4}\\
  E_{43}\mult E_{14} \equiv E_{43}E_{14} .\label{1223r5}
  \end{gather}
Combining the above equalities and taking into account that $[E_{12},E_{23}]=e_{13}\equiv 0$,
we get~\rf{proof12} in~$\Z_4$. We could apply here the stability arguments (as we shall do in the sequel)
but we give some more details at this point to give a f\/lavor of how such derivations of relations work.
For the same, as \rf{1223r1}--\rf{1223r5}, calculations for $\Z_n$, the analogues of the conditions (i)
and (ii), see paragraph 1 of this subsection, restrict $\lambda$ to be of the form $\ve_1-\ve_2+\ve_3-\ve_k$, $k\geq 3$ for
\rf{1223r1}; $\ve_2-\ve_k$, $k\geq 2$ for \rf{1223r2}; $\ve_3-\ve_k$, $k\geq 3$ for~\rf{1223r3}
and \rf{1223r4}; $\ve_4-\ve_k$, $k\geq 4$ for~\rf{1223r5}. It follows, for, say, $n=5$, that the right
hand sides of \rf{1223r1}--\rf{1223r5} might be modif\/ied only by an addition of the term proportional to
$E_{53}E_{15}$; and this will be compensated by an addition of the term, proportional to
$\s_{53}\mult\s_{15}$ to the right hand side of~\rf{proof12}, since $E_{53}\mult E_{15}\equiv E_{53}E_{15}$;
the proportionality coef\/f\/icient is uniquely def\/ined.
This pattern clearly continues and we conclude that there is a relation in $\Z_n$ of the form
\begin{gather}\label{proof13}
\s_{23}\mult\s_{12}-\s_{12}\mult\s_{23}=t_{12}\mult\s_{13} \frac{1}{\th_{12}}+t_{23}\mult\s_{13}\frac{1}{\th_{23}}-
\sum_{k>3}\s_{k3}\mult\s_{1k}X_k  ,
\end{gather}
with certain, uniquely def\/ined, coef\/f\/icients $X_k\in\Uh$, $k=4\lcd n$, and already known $X_4=(\th_{34}+1)
{\th_{34}^{-1}\th_{24}^{-1}}$. To f\/ind $X_5,\dots ,X_n$, we apply to \rf{proof13} the
automorphisms $\q_k$, $k=4\lcd n-1$, which leave invariant the left hand side and the f\/irst two terms in the
right hand side of~\rf{proof13}. The uniqueness of the relation of the form~\rf{proof13}, together with the equality
$\q_k(\s_{k3}\mult\s_{1k})=\s_{k+1,3}\mult\s_{1,k+1}(\th_{k,k+1}+1)\th_{k,k+1}^{-1}$, imply the recurrence relation
$X_{k+1}=\q_k(X_k)\cdot(\th_{k,k+1}+1)\th_{k,k+1}^{-1}$ and we f\/ind
\begin{gather*}%\label{proof14}
X_k=\frac{1}{\th_{2k}}\prod_{j=3}^{k-1}\frac{\th_{jk}+1}{\th_{jk}}  .
\end{gather*}

After the renormalization~\rf{not8} and the change of variables~\rf{3.1a}, the derived relation becomes one of the
relations in the f\/irst line of~\rf{relation3}.


Applying the transformations from the braid group, we obtain the rest of the relations from the list~\rf{relation3}.

{\bf 3b.} We have the following equalities in $\Z_3$:
\begin{gather}\label{proof15}
\s_{12}\mult t_1=t_1\mult\s_{12}  {\ds\frac{\th_{12}+2}{\th_{12}+1}}-t_2\mult\s_{12}  {\ds\frac{1}{\th_{12}+1}}
-\s_{32}\mult\s_{13}  {\ds\frac{\th_{23}+1}{\th_{23}(\th_{13}+1)}}  ,\\
\label{proof16}
\s_{12}\mult t_2=-t_1\s_{12}  {\ds\frac{1}{\th_{12}+1}}+t_2\mult\s_{12}
 {\ds\frac{\th_{12}+2}{\th_{12}+1}}+\s_{32}\mult\s_{13}  {\ds\frac{\th_{13}+2}{\th_{23}(\th_{13}+1)}}  ,
\end{gather}
and the equality in $\Z_4$:
\begin{gather}\label{proof17}
\s_{12}\mult t_4=t_4\mult\s_{12}-\s_{42}\mult\s_{14}  {\ds\frac{\th_{12}+1}{(\th_{14}+1)\th_{24}}}  .
\end{gather}
Equalities \rf{proof15} and \rf{proof16} are the results of the following calculations for $\Z_3$,
using \rf{proof3a}--\rf{proof3c}, and of the commutativity $[H_1,E_{12}]=e_{12}\equiv 0$,
$[H_2,E_{12}]=-e_{12}\equiv 0$:
\begin{gather*}E_{12}\mult H_1 \equiv E_{12}H_1+H_{12}E_{12}  \frac{1}{\th_{12}+1}-
E_{32}E_{13}  \frac{\th_{12}}{(\th_{12}+1)(\th_{13}+1)}  , \\
 E_{12}\mult H_2 \equiv E_{12}H_2-H_{12}E_{12}  \frac{1}{\th_{12}+1}-
E_{32}E_{13}  \frac{1}{(\th_{12}+1)(\th_{13}+1)}, \\
 H_1\mult E_{12} \equiv H_1E_{12}  , \qquad
 H_2\mult E_{12} \equiv H_2E_{12}-E_{32}E_{13}  \frac{1}{\th_{23}} , \qquad
 E_{32}\mult E_{13} \equiv E_{32}E_{13}  .\end{gather*}
The derivation of \rf{proof17} can be done with the help of the following calculations for $\Z_4$:
\begin{gather}\label{pr14n}
E_{12}\mult H_4 \equiv E_{12}H_4+E_{42}E_{14}  \frac{1}{\th_{14}+1},\\
 H_4\mult E_{12} \equiv H_4E_{12}+E_{42}E_{14}  \frac{1}{\th_{24}}  ,\qquad
 E_{42}\mult E_{14} \equiv E_{42}E_{14}  .\nonumber
 \end{gather}
We shall make a comment about the line \rf{pr14n} only. Here one might expect, by  the analogues of the conditions (i) and (ii), see paragraph 1 of this subsection,
non-trivial contributions to $E_{12}\mult H_4$ from the weights 0, $\ve_1-\ve_2$, $\ve_1-\ve_3$ and $\ve_1-\ve_4$.
So we need, in addition to \rf{proof3a}--\rf{proof3c}, some information about
$\J_{\ve_1-\ve_4}$. It follows from the ABRR equation that $\J_{\ve_1-\ve_4}(\th_{14}+1)=-T_{\ve_1-\ve_2}(\J_{\ve_2-\ve_4})
-T_{\ve_1-\ve_3}(\J_{\ve_3-\ve_4})-T_{\ve_1-\ve_4}(\J_0)-T_{\ve_2-\ve_3}(\J_{\ve_1-\ve_2+\ve_3-\ve_4})-
T_{\ve_2-\ve_4}(\J_{\ve_1-\ve_2})-T_{\ve_3-\ve_4}(\J_{\ve_1-\ve_3})$. Since $e_{13}$ and $e_{12}$
commute with $H_4$, the parts $T_{\ve_2-\ve_4}(\J_{\ve_1-\ve_2})$ and $T_{\ve_3-\ve_4}(\J_{\ve_1-\ve_3})$
do not contribute; $e_{42}$ and $e_{43}$ commute with $E_{12}$, so the parts  $T_{\ve_1-\ve_2}(\J_{\ve_2-\ve_4})$
and $T_{\ve_1-\ve_3}(\J_{\ve_3-\ve_4})$ do not contribute either; $\J_{\ve_1-\ve_2+\ve_3-\ve_4}=
\J_{\ve_3-\ve_4}\J_{\ve_1-\ve_2}$ does not contribute again since $e_{12}$ commute with $H_4$.
Thus the only contribution is from $T_{\ve_1-\ve_4}(\J_0)$ and we quickly arrive at~\rf{pr14n}.


Applying the automorphism $\q_3$ of the algebra $\Z_4$ to the relation \rf{proof17}, see \rf{acwot} and~\rf{3.5}, we f\/ind
\begin{gather*}
\s_{12}\mult\big(t_3\th_{34}-t_4\big)=\big(t_3\th_{34}-t_4\big)\mult \s_{12}
-\s_{32}\mult\s_{13}  {\ds\frac{\th_{34}(\th_{12}+1)}{(\th_{13}+1)\th_{23}}}  .
\end{gather*}
We then add \rf{proof17} to this relation  and obtain the following relation in $\Z_4$:
\begin{gather}\label{proof18}
\s_{12}\mult t_3=t_3\mult\s_{12}
-\s_{32}\mult \s_{13}  {\ds\frac{(\th_{12}+1)}{(\th_{13}+1)\th_{23}}}-
\s_{42}\mult\s_{14}  {\ds\frac{\th_{12}+1}{(\th_{14}+1)\th_{24}\th_{34}}}  .
\end{gather}

The stabilization arguments for \rf{proof15}, \rf{proof16} and \rf{proof18} imply the existence of the following relations in $\Z_n$:
\begin{gather}\label{proof19}\s_{12}\mult t_1 =t_1\mult\s_{12}  {\ds\frac{\th_{12}+2}{\th_{12}+1}}-
t_2\mult\s_{12}  {\ds\frac{1}{\th_{12}+1}}+\sum_{k>2}\s_{k2}\mult\s_{1k}X_k^{(1)},\\
 \label{proof20}\s_{12}\mult t_2 =-t_1\mult\s_{12}  {\ds\frac{1}{\th_{12}+1}}+t_2\mult\s_{12}  {\ds\frac{\th_{12}+2}{\th_{12}+1}}
+\sum_{k>2}\s_{k2}\mult\s_{1k}X_k^{(2)},
\\
\label{proof21}\s_{12}\mult t_3 =t_3\mult\s_{12}-\s_{32}\mult\s_{13}  {\ds\frac{(\th_{12}+1)}{(\th_{13}+1)\th_{23}}}
+\sum_{k>3}\s_{k2}\mult\s_{1k}X_k^{(3)},
\end{gather}
where all $X_k^{(i)}$ belong to $\Uh$ and the initial $X_k^{(i)}$ are known:
\begin{gather*} X_3^{(1)} =-\frac{\th_{23}+1}{\th_{23}(\th_{13}+1)}  ,\qquad
X_3^{(2)} =\frac{\th_{13}+2}{\th_{23}(\th_{13}+1)}  ,\qquad
X_4^{(3)} =-\frac{\th_{12}+1}{(\th_{14}+1)\th_{24}\th_{34}}  .
\end{gather*}

By the braid group transformation laws, $X_{k+1}^{(i)}=\q_k\big(X_k^{(i)}\big)\cdot(\th_{k,k+1}+1)\th_{k,k+1}^{-1}$
with $k>2$ for $i=1,2$ and $k>3$ for $i=3$, so that
\begin{gather*}%\label{proof22}
X_k^{(1)}=-\frac{1}{\th_{1k}+1}\prod_{j=2}^{k-1}\frac{\th_{jk}+1}{\th_{jk}} ,\qquad
X_k^{(2)}=\frac{\th_{1k}+2}{(\th_{1k}+1)\th_{2k}}\prod_{j=3}^{k-1}
\frac{\th_{jk}+1}{\th_{jk}}  ,\\
X_k^{(3)}=-\frac{\th_{12}+1}{(\th_{1k}+1)\th_{2k}\th_{3k}}\prod_{j=4}^{k-1}
\frac{\th_{jk}+1}{\th_{jk}}  .
\end{gather*}

After the renormalization \rf{not8} and the change of variables~\rf{3.1a}, the relations \rf{proof19}--\rf{proof21} turn into
the relations~\rf{3.12} for $i=1$, $j=2$ and $k=3$.


Applying the transformations from the braid group, we obtain the rest of the relations from the list \rf{3.12}.

{\bf 4a.} We now prove the relations \rf{relation4a} using the arguments similar to \cite[Subsection 6.1.2]{Zh}.
 Consider the
products $H_k\mult_\lambda H_l$  and $H_l\mult_\lambda H_k$ with $\lambda\not=0$.
These products are linear combinations, over $\Uh$, of monomials
\begin{gather*}
a_{kl;\vec{\gamma}}:=H_ke_{-\gamma_1}\cdots e_{-\gamma_m}e_{\gamma_m}\cdots e_{\gamma_1}H_l\qquad\text{and}\qquad
a_{lk;\vec{\gamma}}:=H_le_{-\gamma_1}\cdots e_{-\gamma_m}e_{\gamma_m}\cdots e_{\gamma_1}H_k  ,
\end{gather*}
respectively; here $m\geq 0$ and $\vec{\gamma}:=\{\gamma_1,\dots ,\gamma_m\}$. By construction, the coef\/f\/icient, from $\Uh$, of the monomial
$a_{kl;\vec{\gamma}}$ in $H_k\mult_\lambda H_l$ equals the coef\/f\/icient of $a_{lk;\vec{\gamma}}$ in $H_l\mult_\lambda H_k$.
The expressions $a_{kl;\vec{\gamma}}$ and $a_{lk;\vec{\gamma}}$ are both equal in $\Z_n$  to
\begin{gather*}(\gamma_1,\ve_k)(\gamma_1,\ve_l)E_{-\gamma_1}e_{-\gamma_2}\cdots
e_{-\gamma_m}e_{\gamma_m}\cdots e_{\gamma_2}E_{\gamma_1}  .\end{gather*}
Thus $H_k\mult_\lambda H_l\equiv H_l\mult_\lambda H_k$  for any $\lambda\not=0$. In the zero weight part $\mult_0$
of the product $\mult$ we have the equality $H_kH_l=H_lH_k$ as well. Therefore,
$H_k\mult H_l\equiv H_l\mult H_k$.

{\bf 4b.} The last group \rf{relation4} of relations is left. Like above, we f\/irst explicitly derive the following relation in~$\Z_3$:
\begin{gather} \s_{12}\mult \s_{21}=  h_{12}-t_{12}\mult t_{12}  \frac{1}{\th_{12}-1}+\s_{21}\mult\s_{12}
\frac{(\th_{12}-1)(\th_{12} +2)}{\th_{12}(\th_{12} +1)}\nonumber\\
\phantom{\s_{12}\mult \s_{21}=}{} +  \s_{31}\mult\s_{13} \frac{(\th_{12}-1) (\th_{13} +2)}{\th_{12}\th_{23}(\th_{13}+1)}
-\s_{32}\mult\s_{23}  \frac{\th_{23} +2}{(\th_{23} +1)\th_{13}}  \label{re01}
\end{gather}
(the f\/irst term in the right hand side is $h_{12}$, without hat). The relation \rf{re01} is a corollary of the following calculations
for $\Z_3$, together with the commutation relation $[E_{12},E_{21}]=h_{12}$,
\begin{gather}E_{12}\mult E_{21}\equiv   E_{12}E_{21}-H_{12}^2 \frac{1}{(\th_{12}-1)}
+E_{21}E_{12} \frac{2}{(\th_{12}-1)\th_{12}}\nonumber\\
\phantom{E_{12}\mult E_{21}\equiv}{} - E_{32}E_{23} \frac{\th_{12} -2}{(\th_{12}-1)\th_{13}}+
E_{31}E_{13} \frac{\th_{12} -2}{(\th_{12}-1)\th_{12}\th_{13}} ,\label{zemod1}\\
 H_{12}\mult H_{12}\equiv   H_{12}^2-E_{21}E_{12} \frac{4}{\th_{12} +1}
-E_{32}E_{23} \frac{1}{\th_{23} +1}\nonumber\\
\phantom{H_{12}\mult H_{12}\equiv}{} +  E_{31}E_{13}\left(-1+\frac{1}{\th_{23} +1}+\frac{4}{\th_{12} +1}\right)\frac{1}{\th_{13}+1} ,\label{zemod2}\\
E_{21}\mult E_{12}\equiv E_{21}E_{12}-E_{31}E_{13} \frac{1}{\th_{23}} ,\label{zemod3}\\
 E_{32}\mult E_{23}\equiv  E_{32}E_{23}-E_{31}E_{13} \frac{1}{\th_{12}} ,\label{zemod4}\\
 E_{31}\mult E_{13}\equiv  E_{31}E_{13} .\label{zemod5}
 \end{gather}
Here only the calculation of $E_{12}\mult E_{21}$ deserves a little explanation; by  the analogues of the conditions (i) and (ii),
see paragraph 1 of this subsection, non-trivial contributions to $E_{12}\mult E_{21}$ from the weights 0,
$\ve_1-\ve_2$, $2(\ve_1-\ve_2)$, $\ve_1-\ve_3$ and $2\ve_1-\ve_2-\ve_3$ are possible. By the ABRR equation,
$\J_{2(\ve_1-\ve_2)}(2\th_{12}+4)=-T_{\ve_1-\ve_2}(\J_{\ve_1-\ve_2})$ and $\J_{2\ve_1-\ve_2-\ve_3}
(2\th_1-\th_2-\th_3+3)=-T_{\ve_1-\ve_2}(\J_{\ve_1-\ve_3})-T_{\ve_1-\ve_3}(\J_{\ve_1-\ve_2})-T_{\ve_2-\ve_3}
(\J_{2(\ve_1-\ve_2)})$. We leave further details to the reader.


By the stabilization law in $\Z_4$ we have a relation
\begin{gather} \label{proof23}
\s_{12}\mult \s_{21}= h_{12}-t_{12}\mult t_{12}  \frac{1}{\th_{12}-1}
+\sum_{1\leq i<j\leq n}\s_{ji}\mult\s_{ij}X_{ij} ,\qquad X_{ij}\in\Uh
\end{gather}
with $n=4$, which dif\/fers from \rf{re01} by a presence of terms
\begin{gather*}\s_{43}\mult\s_{34} ,\qquad \s_{42}\mult\s_{24} ,\qquad \s_{41}\mult\s_{14} ,
\end{gather*}
with coef\/f\/icients in $\Uh$. Consider in $\Z_4$ the products  $\s_{12}\mult\s_{21}$, $t_{12}\mult t_{12}$
and $\s_{ji}\mult\s_{ij}$, $1\leq i<j\leq 4$. The weights $(\ve_3-\ve_4)-(\ve_i-\ve_j)$ do not belong
to the cone $\Q_+$ if $1\leq i<j<4$. Thus in the decomposition
\begin{gather*}
E_{ji}\mult E_{ij}\equiv \sum_{k<l}E_{lk}E_{kl}a_{kl} ,\qquad a_{kl}\in\Uh ,\quad 1\leq i<j<4 ,
\end{gather*}
the term with $E_{43}E_{34}$ has a zero coef\/f\/icient, $a_{34}=0$. The same statement holds for the products
$E_{41}\mult E_{14}$ and $E_{42}\mult E_{24}$ since the weights $(\ve_3-\ve_4)-(\ve_i-\ve_4)$ do
not belong to $\Q_+$ for $i=1,2$. Consider the product $E_{12}\mult E_{21}$. Here the term with
$E_{43}E_{34}$ is equal to $E_{12}\mult_{\ve_1-\ve_2+\ve_3-\ve_4}E_{21}$. By \rf{proof3c},
$\J_{\ve_1-\ve_2+\ve_3-\ve_4}=e_{43}e_{21}\otimes e_{12}e_{34}(\th_{12}+1)^{-1}(\th_{34}+1)^{-1}$ and
\begin{gather*}E_{12}\mult_{\ve_1-\ve_2+\ve_3-\ve_4}E_{21}=
E_{12}e_{43}e_{21}e_{12}e_{34}E_{21} \frac{1}{(\th_{12}-1)(\th_{34}+1)}\equiv 0,\end{gather*}
since $[e_{34},E_{21}]=0$ (and $[E_{12},e_{43}]=0$). In the similar manner, the term with
$E_{43}E_{34}$ in $H_{12}\mult H_{12}$ equals $H_{12}\mult_{\ve_3-\ve_4}H_{12}$ and vanishes
since $[e_{34},H_{12}]=0$.

On the other hand, the product $E_{43}\mult E_{34}$
def\/initely contains $E_{43}E_{34}=E_{43}\mult_0 E_{34}$. We thus conclude that the term
$\s_{43}\mult \s_{34}$ is absent in~\rf{proof23}, that is $X_{34}=0$.


For $n>4$, again by the stabilization law, we have a unique relation of the form \rf{proof23}. By uniqueness, it is invariant with respect to the
transformations $\q_3,\q_4,\dots ,\q_{n-1}$ which do not change the product $\s_{12}\mult \s_{21}$. Since $X_{34}=0$, we f\/ind, applying
$\q_4, \q_5, \dots ,\q_{n-1}$,  that  $X_{3j}=0$, $j>3$, wherefrom we further conclude, applying $\q_3,\q_4,\dots ,\q_{j-2}$, that
$X_{ij}=0$, $2<i<j$. We get f\/inally the following relation in $\Z_n$:
\begin{gather} \label{soo1221}
\s_{12}\mult \s_{21}=h_{12}-t_{12}\mult t_{12} \frac{1}{\th_{12}-1}
+\sum_{k=2\lcd n}\s_{k1}\mult\s_{1k}X_{1k}+\sum_{k=3\lcd n}\s_{k2}\mult\s_{2k}X_{2k}
\end{gather}
with known
\begin{gather*}
X_{12}=\frac{(\th_{12}-1)(\th_{12} +2)}{\th_{12}(\th_{12} +1)} ,\qquad
X_{13}=\frac{(\th_{12}-1) (\th_{13} +2)}{\th_{12}\th_{23}(\th_{13}+1)} ,\qquad
X_{23}=-\frac{\th_{23} +2}{(\th_{23} +1)\th_{13}} .
\end{gather*}
Applying to \rf{soo1221} the transformations $\q_3, \q_4,\dots ,\q_{n-1}$ we f\/ind by uniqueness
\begin{gather*}X_{i,k+1}=\frac{\th_{k,k+1}+1}{\th_{k,k+1}} \cdot\q_k(X_{ik}) ,\qquad i=1,2 ;\quad k=3,4,\dots ,n-1 ,
\end{gather*}
and thus
\begin{gather*}
 X_{1k}=\frac{(\th_{12}-1) (\th_{1k} +2)}{\th_{12}\th_{2k}(\th_{1k}+1)}\cdot\prod_{a=3}^{k-1}\frac{\th_{ak}+1}{\th_{ak}}  ,
\qquad X_{2k}=-\frac{\th_{2k} +2}{(\th_{2k} +1)\th_{1k}}\cdot\prod_{a=3}^{k-1}\frac{\th_{ak}+1}{\th_{ak}}
\end{gather*}
for $k>2$.


After the renormalization \rf{not8} and the change of variables~\rf{3.1a}, the relations~\rf{soo1221}  with the
obtained~$X_{1k}$ and $X_{2k}$ turns into the relation~\rf{relation4} for $i=1$ and $j=2$.


Applying the transformations from the braid group, we obtain the rest of the relations from the list~\rf{relation4}.

\subsection{Proof of Theorem \ref{mthe}}\label{proofmain}

For the proof of Theorem \ref{mthe} we just apply the results of~\cite{KO3}, which state that the system ${\mathfrak R}$ is the system of
def\/ining relations once it is satisf\/ied in the algebra $\Z_n$.

\begin{remark}
An attentive look shows that the system $\mathfrak{R}$ is closed under the anti-involution~$\epsilon$; that is, $\epsilon$~transforms any
relation from $\mathfrak{R}$ into a linear over $\Uh$ combination of relations from~$\mathfrak{R}$. Moreover, $\mathfrak{R}$ and
$\epsilon (\mathfrak{R})$ are equivalent over $\Uh$. Indeed, all relations in Section~\rf{sectionDerRel} were derived in
three steps: f\/irst we derive a relation in $\Z_n$ with $n\leq 4$; next by the stabilization principle we extend the derived relation
to~$\Z_n$ with arbitrary~$n$;
and then we f\/ind the whole list of relations of a given (sub)type by applying the braid group transformations (products of the
generators $\q_i$). Due to \rf{invr} one could use $\q_i^{-1}$ instead of $\q_i$ equivalently over
$\Uh$. A~straightforward calculation establishes the equivalence of the extended to arbitrary $n$ lists $\mathfrak{R}$ and
$\epsilon (\mathfrak{R})$ over $\Uh$ for $\Z_n$ with
$n\leq 4$ (this verif\/ication is lengthy for some relations).
Then with the help of~\rf{qeps} we f\/inish the check of the equivalence of $\mathfrak{R}$ and $\epsilon (\mathfrak{R})$ over $\Uh$ for $\Z_n$ with
arbitrary $n$.


Similar arguments establish  the equivalence of $\mathfrak{R}$ and $\omega (\mathfrak{R})$ over $\Uh$; here $\omega$ is the involution def\/ined
in~\rf{not2a}. In~\cite{KO3} this equivalence was obtained dif\/ferently, as a by-product of the equivalence, over $\Uh$, of the system $\mathfrak{R}$ and
the system  \rf{not6} of ordering relations.
\end{remark}
